\section{Devil in the Details}

\begin{quotex}
All religion and all wisdom is reducible, extrinsically and from the human standpoint, to these four laws:

\begin{itemize}


\item enshrined in every tradition is to be observed an Immutable Truth

\item a law of attachment to the Real

\item remembrance of love of God

\item prohibitions and injunctions
\end{itemize}

These make up a fabric of elementary certainties which encompasses and resolves human uncertainty, and thus reduces the whole problem of earthly existence to a geometry that is at once simple and primordial.
\flright{\textsc{Frithjof Schuon}, \emph{Logic and Transcendence}}
\end{quotex}

Here we see Schuon describing the initiatic path as ``simple and primordial''. It has been well marked out by those who know. \textbf{Rene Guenon} likewise points out that the path is precise and exact. Both Schuon and Guenon make use of geometrical symbolism, to indicate the order behind Traditional thinking. Why do you suppose that Plato refused his teaching to anyone ignorant of geometry? There is no room here for disputation or disagreement. It is understood or it is not. Schuon writes:

\begin{quotex}
This principal polarization is refracted endlessly in the universe, but it does so in an unequal manner --- according to the exigencies of manifesting Possibility --- and the subjectivies as a result are not epistemologically equivalent.
\end{quotex}


This explains the disorder in the world, since as the Possibilities manifest, men have different intellectual capacities. Yet there must be some basic, or latent, intelligence to even grasp that. This is all the more true in the modern world when every man believes he is equal and his opinion is as good as another; this is even true, unfortunately, of those who pay lip service to tradition and anti-egalitarianism. In practice, however, it is much more difficult to swallow.

The first of the \textbf{corporal works of mercy} is ``to feed the hungry''. The hungry never object to being fed. But the corresponding \textbf{spiritual work of mercy} is ``to instruct the ignorant''. However, the ignorant, or spiritually hungry, are typically indignant about this, unlike those with physical hunger. Those who sincerely desire true knowledge lap it up, but they are unequally distributed. Being overwhelmed by the world is a sure way to remain in ignorance, and a reorientation is necessary. Schuon describes it this way:

\begin{quotex}
Man has around him the bewildering multitude of phenomena, perfect intelligence will consist in perceiving their homogeneity and their outwardness as proceeding respectively from a transcendent unity and a reality that is inward; the world will then appear, not as an incoherent mass of quasi-absolute phenomena, but as a single veil into which the phenomena are woven; in this veil, they are joined but not confused distinct but not separated. In the center resides the discerning and unifying intelligence, namely, the intelligence that is conscious of the Principle; it is thanks to this consciousness alone that the phenomenal world can appear both in its substantial homogeneity and in its contingency, its outwardness, its nothingness.
\end{quotex}


In \textit{East and West}, Guenon explains:

\begin{quotex}
What we call a normal civilization is one that is based on principles, in the true sense of this word, one where everything is arranged in hierarchy to conform to these principles, so that everything in it is seen as the application and extension of a doctrine whose essence is purely intellectual and metaphysical: this is what we mean also when we speak of a traditional civilization
\end{quotex}


So, again what is required is knowledge of principle, not knowledge of the world. As one of the superstitions of life, Guenon makes this clear; the opposite of the knowledge of principles is an attachment to the sensate world. This can be overcome only by ``an understanding of the universal order, which is beyond the reach of the modern world where disorder is law.'' Guenon Unfortunately, ``the modern mind faces almost exclusively outward, towards the world of the senses.'' To expect order to reveal itself in the sensate world is misguided and delusional; it will lead nowhere. To achieve anything in this area requires efforts and fortitude, as well as both the desire and the ability to turn away from it.

Modern man is also sentimental, driven by the inducements of the emotions. He loves florid language and dreams his Walter Mitty\footnote{\url{https://en.wikipedia.org/wiki/Walter\_Mitty}} dreams of being a warrior, high initiate or sage. Far be it from us to discourage such noble goals, but, alas, the devil is in the details. It is one thing to talk about overcoming ``lower powers'' and the like; it is quite another to actually name them. Of course, the battle against such base powers needs to be fought within, long before they can be fought without.

We can get instructions for this battle from various sources, such as, \emph{The Yoga Sutras of Patanjali} or the eight-fold path described by Buddha. Closer to home, we have the \emph{Divine Comedy} by Dante, known to be an initiate of the order of the \emph{Fedeli d'Amore}. In great detail, he describes the obstacles that keep us bound to the horrors of the world, the steps required to purge ourselves of the influences of those obstacles, and finally the final stage to the Supreme Identity.

To be clear, and as Guenon insists, this is not some holy mysticism or wishful thinking that it the result of a divine accident. Rather, it is a logical and coherent path that can be followed by those who understand it and make the efforts to follow it. It is easy to talk about ``liberation'' in some abstract sense. But we are, or should be, interested in understanding specifically what that means. Rene Guenon has made a precise definition of that end result\footnote{\url{https://gornahoor.net/?p=4527}}, so it can be understood by anyone of intelligence willing to make the effort. There are some men who are no longer satisfied by vague ideas and an attachment to the contemporary world. They will be fed.


\flrightit{Posted on 2012-07-31 by Cologero}

